\chapter{The Pipe Operator}
\label{cap:pipes}

\section{Basic syntax}

The \lstinline{|>} operator passes the value on its left as the first
argument to the function on its right:

\begin{lstlisting}
expr |> f(args...)
// equivalent to: f(expr, args...)
\end{lstlisting}

Without pipe:
\begin{lstlisting}
display sort(filter(df, income > 0), income desc)
\end{lstlisting}

With pipe:
\begin{lstlisting}
df |> filter(income > 0) |> sort(income desc) |> display
\end{lstlisting}

The pipe allows reading transformations in the order they occur, left to
right (or top to bottom), without nested parentheses.

\section{Semantics with DataFrames}

The pipe's behaviour with DataFrames follows a single rule:

\begin{center}
\begin{tabular}{ll}
\toprule
\textbf{Form} & \textbf{Effect} \\
\midrule
\texttt{df |> f(...)}           & Assign-back: df receives the result \\
\texttt{let df2 = df |> f(...)} & Functional: df unchanged, df2 is new \\
\bottomrule
\end{tabular}
\end{center}

\subsection*{Assign-back}

When the pipe starts with a variable (\lstinline{df |> ...}) and there is
no \lstinline{let}, the result is assigned back to the source variable:

\begin{lstlisting}[caption={Assign-back: df is modified}]
load "data.csv" as df

df |> filter(year >= 2010)
   |> dropna
   |> sort(income desc)

// df is now the result of all transformations
\end{lstlisting}

\subsection*{Functional form with \texttt{let}}

\begin{lstlisting}[caption={let preserves the original}]
load "data.csv" as df

let df_clean = df
  |> filter(year >= 2010)
  |> dropna

// df is intact; df_clean is the filtered version
\end{lstlisting}

\section{COW and pipes}

COW semantics ensures that assign-back is efficient: while \lstinline{df}
and \lstinline{df_clean} share the same data (before any mutation), no
memory copy occurs. The actual copy only happens when the data diverges.

\begin{lstlisting}[caption={Preserving baseline with let}]
load "panel.csv" as df
let baseline = df      // cheap alias

df |> filter(treated == 1)   // modifies df, not baseline

display count(baseline)    // all original obs
display count(df)          // only the treatment group
\end{lstlisting}

\section{Long chains}

The pipe is especially useful in data cleaning and transformation
pipelines:

\begin{lstlisting}[caption={Complete preparation pipeline}]
load "survey.csv" as survey

let sample = survey
  |> filter(age >= 18 and age <= 65)
  |> filter(hours_worked > 0)
  |> dropna
  |> generate(lincome = log(income + 1))
  |> generate(exp  = age - years_educ - 6)
  |> generate(exp2 = exp^2)
  |> sort(state year)

display count(sample)
\end{lstlisting}

\section{Using placeholders}

If you need to pass the left-hand side value to a specific argument position other than the first one, you can use the underscore (\lstinline{_}) as a placeholder:

\begin{lstlisting}
expr |> f(a, _)
// equivalent to: f(a, expr)
\end{lstlisting}

This is particularly useful when piping DataFrames into estimators (like \lstinline{ols} or \lstinline{iv}), where the DataFrame is not the first argument:

\begin{lstlisting}
let model = df |> ols(y ~ x1 + x2, _)
\end{lstlisting}

\section{Pipe with closures}

The \lstinline{|>} operator also works with closures:

\begin{lstlisting}
let result = 5 |> fn(x) { x^2 + 1 }
display result    // 26
\end{lstlisting}
